\documentclass[12pt,a4paper]{article}
\usepackage[utf8]{inputenc}
\usepackage{amsmath}
\usepackage{amssymb}
\usepackage{geometry}
\usepackage{hyperref}
\usepackage{bm}

\geometry{margin=2.5cm}

\title{Mathematical Formulas - Speech Enhancement Project}
\author{}
\date{}

\begin{document}

\maketitle

\tableofcontents
\newpage

\section{Signal Processing Fundamentals}

\subsection{Short-Time Fourier Transform (STFT)}

The STFT transforms a time-domain signal into the time-frequency domain:

\begin{equation}
X(f, t) = \int_{-\infty}^{\infty} x(\tau) \cdot w(\tau - t) \cdot e^{-j2\pi f\tau} d\tau
\end{equation}

Discrete STFT for frame $m$:
\begin{equation}
X[k, m] = \sum_{n=0}^{N-1} x[n + mH] \cdot w[n] \cdot e^{-j2\pi kn/N}
\end{equation}

where:
\begin{itemize}
    \item $x[n]$ is the input signal
    \item $w[n]$ is the window function (Hanning window)
    \item $N$ is the FFT size (n\_fft = 512)
    \item $H$ is the hop length (hop\_length = 160)
    \item $k$ is the frequency bin index
    \item $m$ is the frame index
\end{itemize}

\subsection{Inverse STFT (ISTFT)}

Reconstruction from STFT:
\begin{equation}
x[n] = \frac{\sum_{m} w[n - mH] \cdot \text{IFFT}\{X[k, m]\}}{\sum_{m} w^2[n - mH]}
\end{equation}

\subsection{Magnitude and Phase}

From complex STFT:
\begin{align}
\text{Magnitude:} \quad |X[k, m]| &= \sqrt{\text{Re}(X[k,m])^2 + \text{Im}(X[k,m])^2} \\
\text{Phase:} \quad \angle X[k, m] &= \arctan\left(\frac{\text{Im}(X[k,m])}{\text{Re}(X[k,m])}\right)
\end{align}

Reconstruction:
\begin{equation}
X[k, m] = |X[k, m]| \cdot e^{j\angle X[k, m]}
\end{equation}

\section{Speech Enhancement Methods}

\subsection{Ideal Ratio Mask (IRM)}

Target mask for neural network training:
\begin{equation}
\text{IRM}[k, m] = \frac{|S_{\text{clean}}[k, m]|}{|S_{\text{clean}}[k, m]| + |N[k, m]| + \epsilon}
\end{equation}

where:
\begin{itemize}
    \item $S_{\text{clean}}$ is the clean speech STFT
    \item $N$ is the noise STFT (approximated as $S_{\text{noisy}} - S_{\text{clean}}$)
    \item $\epsilon = 10^{-12}$ prevents division by zero
\end{itemize}

Enhanced speech:
\begin{equation}
\hat{S}_{\text{clean}}[k, m] = \text{IRM}[k, m] \cdot S_{\text{noisy}}[k, m]
\end{equation}

\subsection{Spectral Subtraction}

Basic spectral subtraction:
\begin{equation}
|\hat{S}[k, m]| = \max\left(|Y[k, m]| - \alpha \cdot |\hat{N}[k]|, \beta \cdot |\hat{N}[k]|\right)
\end{equation}

where:
\begin{itemize}
    \item $Y[k, m]$ is the noisy speech magnitude
    \item $\hat{N}[k]$ is the estimated noise magnitude (averaged over silence frames)
    \item $\alpha$ is the over-subtraction factor (default: 1.0)
    \item $\beta$ is the spectral floor factor (default: 0.02)
\end{itemize}

Adaptive noise tracking:
\begin{equation}
\hat{N}[k, t] = \lambda \cdot \hat{N}[k, t-1] + (1 - \lambda) \cdot |Y[k, t]|
\end{equation}

where $\lambda = 0.95$ is the smoothing factor (applied only during silence frames).

\subsection{Wiener Filter}

Wiener gain function:
\begin{equation}
G[k, m] = \frac{\text{SNR}_{\text{post}}[k, m]}{\text{SNR}_{\text{post}}[k, m] + 1}
\end{equation}

A posteriori SNR:
\begin{equation}
\text{SNR}_{\text{post}}[k, m] = \frac{|Y[k, m]|^2}{\hat{\sigma}_N^2[k, m] + \epsilon}
\end{equation}

Enhanced magnitude:
\begin{equation}
|\hat{S}[k, m]| = G[k, m] \cdot |Y[k, m]|
\end{equation}

Noise PSD estimation (during silence):
\begin{equation}
\hat{\sigma}_N^2[k] = \frac{1}{M} \sum_{m \in \text{silence}} |Y[k, m]|^2
\end{equation}

Adaptive noise PSD tracking:
\begin{equation}
\hat{\sigma}_N^2[k, t] = \lambda \cdot \hat{\sigma}_N^2[k, t-1] + (1 - \lambda) \cdot |Y[k, t]|^2
\end{equation}

\section{Voice Activity Detection (VAD)}

\subsection{Energy-based Detection}

Frame energy:
\begin{equation}
E[m] = \frac{1}{L} \sum_{n=0}^{L-1} x^2[n + mH]
\end{equation}

Normalized energy:
\begin{equation}
E_{\text{norm}}[m] = \frac{E[m]}{\max(E) + \epsilon}
\end{equation}

Speech detection:
\begin{equation}
\text{VAD}_E[m] = \begin{cases}
1 & \text{if } E_{\text{norm}}[m] > \theta_E \\
0 & \text{otherwise}
\end{cases}
\end{equation}

where $\theta_E = 0.1$ is the energy threshold ratio.

\subsection{Zero-Crossing Rate (ZCR)}

ZCR for frame $m$:
\begin{equation}
\text{ZCR}[m] = \frac{1}{L-1} \sum_{n=0}^{L-2} \mathbb{1}\{\text{sign}(x[n + mH]) \neq \text{sign}(x[n + mH + 1])\}
\end{equation}

Speech typically has lower ZCR than noise:
\begin{equation}
\text{VAD}_{\text{ZCR}}[m] = \begin{cases}
1 & \text{if } \text{ZCR}[m] < \theta_{\text{ZCR}} \\
0 & \text{otherwise}
\end{cases}
\end{equation}

where $\theta_{\text{ZCR}} = 0.2$ is the maximum ZCR for speech.

\subsection{Combined VAD}

Final VAD decision:
\begin{equation}
\text{VAD}[m] = \text{VAD}_E[m] \land \text{VAD}_{\text{ZCR}}[m]
\end{equation}

Post-processing to remove short segments:
\begin{itemize}
    \item Remove speech islands shorter than 3 frames
    \item Remove silence islands shorter than 3 frames
\end{itemize}

\section{Evaluation Metrics}

\subsection{Signal-to-Noise Ratio (SNR)}

Global SNR in dB:
\begin{equation}
\text{SNR} = 10 \log_{10} \left(\frac{P_{\text{signal}}}{P_{\text{noise}}}\right)
\end{equation}

where:
\begin{align}
P_{\text{signal}} &= \frac{1}{N} \sum_{n=1}^{N} s^2[n] \\
P_{\text{noise}} &= \frac{1}{N} \sum_{n=1}^{N} (s[n] - \hat{s}[n])^2
\end{align}

\subsection{Segmental SNR}

Frame-based SNR with clipping:
\begin{equation}
\text{SegSNR} = \frac{1}{M} \sum_{m=1}^{M} \text{clip}\left(10 \log_{10} \left(\frac{P_{\text{signal}}^{(m)}}{P_{\text{noise}}^{(m)}}\right), -10, 35\right)
\end{equation}

\subsection{Scale-Invariant SDR (SI-SDR)}

Scale-invariant projection:
\begin{equation}
\alpha = \frac{\langle \hat{s}, s \rangle}{\langle s, s \rangle}
\end{equation}

Target signal:
\begin{equation}
s_{\text{target}} = \alpha \cdot s
\end{equation}

Error signal:
\begin{equation}
e_{\text{noise}} = \hat{s} - s_{\text{target}}
\end{equation}

SI-SDR in dB:
\begin{equation}
\text{SI-SDR} = 10 \log_{10} \left(\frac{\|s_{\text{target}}\|^2}{\|e_{\text{noise}}\|^2 + \epsilon}\right)
\end{equation}

\subsection{Log-Spectral Distance (LSD)}

Per-frame LSD:
\begin{equation}
\text{LSD}[m] = \sqrt{\frac{1}{K} \sum_{k=1}^{K} \left(10\log_{10}P_{\text{clean}}[k,m] - 10\log_{10}P_{\text{proc}}[k,m]\right)^2}
\end{equation}

Average LSD:
\begin{equation}
\text{LSD} = \frac{1}{M} \sum_{m=1}^{M} \text{LSD}[m]
\end{equation}

\subsection{VAD Metrics}

For binary classification (speech vs silence):

True Positive (TP), True Negative (TN), False Positive (FP), False Negative (FN)

\begin{align}
\text{Accuracy} &= \frac{\text{TP} + \text{TN}}{\text{TP} + \text{TN} + \text{FP} + \text{FN}} \\
\text{Precision} &= \frac{\text{TP}}{\text{TP} + \text{FP}} \\
\text{Recall} &= \frac{\text{TP}}{\text{TP} + \text{FN}} \\
\text{F1-Score} &= \frac{2 \cdot \text{Precision} \cdot \text{Recall}}{\text{Precision} + \text{Recall}}
\end{align}

Miss Rate and False Alarm Rate:
\begin{align}
\text{Miss Rate} &= \frac{\text{FN}}{\text{TP} + \text{FN}} = 1 - \text{Recall} \\
\text{False Alarm Rate} &= \frac{\text{FP}}{\text{FP} + \text{TN}}
\end{align}

\section{Audio Normalization}

\subsection{Peak Normalization}

\begin{equation}
x_{\text{norm}}[n] = \frac{x[n]}{\max_n |x[n]|}
\end{equation}

\subsection{RMS Normalization}

Root Mean Square:
\begin{equation}
\text{RMS} = \sqrt{\frac{1}{N} \sum_{n=1}^{N} x^2[n]}
\end{equation}

Current level in dB:
\begin{equation}
L_{\text{current}} = 20 \log_{10}(\text{RMS})
\end{equation}

Gain to reach target level:
\begin{equation}
g_{\text{dB}} = L_{\text{target}} - L_{\text{current}}
\end{equation}

Linear gain:
\begin{equation}
g = 10^{g_{\text{dB}}/20}
\end{equation}

Normalized signal:
\begin{equation}
x_{\text{norm}}[n] = g \cdot x[n]
\end{equation}

\subsection{LUFS Normalization}

ITU-R BS.1770 loudness measurement:

Mean square with K-weighting:
\begin{equation}
\overline{x^2} = \frac{1}{N} \sum_{n=1}^{N} (K_{\text{weight}} \star x[n])^2
\end{equation}

LUFS (Loudness Units relative to Full Scale):
\begin{equation}
\text{LUFS} = -0.691 + 10 \log_{10}(\overline{x^2})
\end{equation}

\subsection{Dynamic Range}

Peak-to-RMS ratio in dB:
\begin{equation}
\text{DR}_{\text{dB}} = 20 \log_{10}\left(\frac{\max_n |x[n]|}{\text{RMS}}\right)
\end{equation}

\section{Data Augmentation}

\subsection{SNR-based Mixing}

Target SNR in linear scale:
\begin{equation}
\rho = 10^{-\text{SNR}_{\text{dB}}/10}
\end{equation}

Noise scaling factor:
\begin{equation}
\alpha = \sqrt{\rho \cdot \frac{P_{\text{clean}}}{P_{\text{noise}}}}
\end{equation}

Mixed signal:
\begin{equation}
y[n] = s[n] + \alpha \cdot n[n]
\end{equation}

\subsection{Gain Adjustment}

Gain in dB to linear:
\begin{equation}
g = 10^{g_{\text{dB}}/20}
\end{equation}

Applied signal:
\begin{equation}
x_{\text{aug}}[n] = g \cdot x[n]
\end{equation}

\subsection{Pitch Shift}

Semitones to rate conversion:
\begin{equation}
r = 2^{s/12}
\end{equation}

where $s$ is the number of semitones (positive for upshift, negative for downshift).

\subsection{Compression}

Compression with attack/release envelope:
\begin{equation}
x_{\text{comp}}[n] = \text{sign}(x[n]) \cdot |x[n]|^{1 - 1/R}
\end{equation}

where $R$ is the compression ratio (e.g., $R=4$ means 4:1 compression).

Threshold-based compression:
\begin{equation}
x_{\text{comp}}[n] = \begin{cases}
x[n] & \text{if } |x[n]| < \theta \\
\text{sign}(x[n]) \cdot \left(\theta + \frac{|x[n]| - \theta}{R}\right) & \text{otherwise}
\end{cases}
\end{equation}

\section{Neural Network Architecture}

\subsection{U-Net Encoder-Decoder}

Encoder (downsampling):
\begin{equation}
h_i = f_{\text{down}}(h_{i-1}), \quad i = 1, 2, 3, 4
\end{equation}

Bottleneck with LSTM:
\begin{equation}
h_{\text{lstm}} = \text{BiLSTM}(h_4)
\end{equation}

Decoder (upsampling) with skip connections:
\begin{equation}
h_i = f_{\text{up}}([h_{i+1}, h_{\text{skip}_i}]), \quad i = 4, 3, 2, 1
\end{equation}

where $[\cdot, \cdot]$ denotes concatenation.

\subsection{Channel Attention}

Attention weights:
\begin{equation}
\alpha_c = \frac{\exp(W_c \cdot \text{GlobalAvgPool}(h_c))}{\sum_{c'} \exp(W_{c'} \cdot \text{GlobalAvgPool}(h_{c'}))}
\end{equation}

Attended features:
\begin{equation}
h'_c = \alpha_c \cdot h_c
\end{equation}

\subsection{Loss Function}

Mean Squared Error (MSE) between predicted and ideal masks:
\begin{equation}
\mathcal{L}_{\text{MSE}} = \frac{1}{KM} \sum_{k=1}^{K} \sum_{m=1}^{M} (\text{IRM}_{\text{pred}}[k,m] - \text{IRM}_{\text{true}}[k,m])^2
\end{equation}

Mean Absolute Error (MAE):
\begin{equation}
\mathcal{L}_{\text{MAE}} = \frac{1}{KM} \sum_{k=1}^{K} \sum_{m=1}^{M} |\text{IRM}_{\text{pred}}[k,m] - \text{IRM}_{\text{true}}[k,m]|
\end{equation}

Combined loss:
\begin{equation}
\mathcal{L} = \alpha \cdot \mathcal{L}_{\text{MSE}} + (1-\alpha) \cdot \mathcal{L}_{\text{MAE}}
\end{equation}

\section{Notation Summary}

\begin{tabular}{ll}
\hline
Symbol & Description \\
\hline
$x[n]$ & Time-domain signal \\
$X[k,m]$ & STFT (frequency $k$, frame $m$) \\
$|X|$ & Magnitude spectrum \\
$\angle X$ & Phase spectrum \\
$s[n]$ & Clean speech signal \\
$n[n]$ & Noise signal \\
$y[n]$ & Noisy signal ($s + n$) \\
$\hat{s}[n]$ & Enhanced/estimated signal \\
$N$ & FFT size (512) \\
$H$ & Hop length (160 samples) \\
$L$ & Window/frame length (400 samples) \\
$\epsilon$ & Small constant ($10^{-12}$) \\
$\lambda$ & Smoothing factor (0.95) \\
$\alpha$ & Over-subtraction factor / weighting \\
$\beta$ & Spectral floor factor (0.02) \\
\hline
\end{tabular}

\section{Implementation Notes}

\subsection{STFT Parameters}
\begin{itemize}
    \item Sample rate: 16000 Hz
    \item FFT size (n\_fft): 512
    \item Hop length: 160 samples (10 ms)
    \item Window length: 400 samples (25 ms)
    \item Window type: Hanning
    \item Frequency bins: 257 (n\_fft/2 + 1)
\end{itemize}

\subsection{VAD Parameters}
\begin{itemize}
    \item Frame length: 400 samples (25 ms)
    \item Hop length: 160 samples (10 ms)
    \item Energy threshold ratio: 0.1
    \item Max ZCR for speech: 0.2
    \item Min speech frames: 3
    \item Min silence frames: 3
\end{itemize}

\subsection{Training Configuration}
\begin{itemize}
    \item Batch size: 8
    \item Optimizer: Adam ($\beta_1=0.9$, $\beta_2=0.999$)
    \item Learning rate: $10^{-3}$
    \item Weight decay: $10^{-5}$
    \item Loss: MSE (Mean Squared Error)
    \item Mixed precision: Enabled (FP16)
\end{itemize}

\end{document}
